\section{The Fracture Gap and Defect Continuity}\label{sec:fracture}

The defect framework proves the universal upper bound
$D(V')\le \Eseq(|V'|)$.  The following two results sharpen it at powers of
two and immediately below powers of two.  They concern root defect, not the
number of internal arrangement-graph edges; arrangement lines are cliques,
so a triangle can have three internal edges while having defect $2$.

\begin{definition}[Coordinate box]\label{def:box}
A $d$-box is a set obtained by choosing $d$ coordinates with two allowed
symbols in each, fixing all remaining coordinates, and taking every injective
combination of those choices.  Thus a $d$-box has $2^d$ vertices.
\end{definition}

\begin{lemma}[Box rigidity]\label{lem:box-rigidity}
If $B$ is a $d$-box and $w\notin B$, then $w$ shares at most one coordinate
root with $B$.  If $B$ and $B'$ are distinct $d$-boxes, then
$|B\cap B'|\le 2^{d-1}$.
\end{lemma}

\begin{lemma}[Power-of-two partition gap]\label{lem:fracture-arith}
Let $h=2^{d-1}$.  For every non-balanced partition of $2^d$,

\[
\sum_i \Eseq(a_i)+2^d-\max_i a_i
 \le \Eseq(2^d)-(d-1).
\]
For $1\le a<h$ the two-part deficit is exactly

\[
\Eseq(2^d)-\Eseq(a)-\Eseq(2^d-a)-a
 =a(d-1)-2\Eseq(a)\ge d-1.
\]
\end{lemma}

\begin{theorem}[Power-of-two fracture gap and equality]\label{thm:fracture}
Let $S\subseteq V(A(n,k))$ with $|S|=2^d$.  Then

\[
D(S)=\Eseq(2^d)
\quad\Longleftrightarrow\quad
S\text{ is a }d\text{-box}.
\]

If $S$ is not a $d$-box, then

\[
D(S)\le \Eseq(2^d)-(d-1).
\]
Thus the intermediate values
$\Eseq(2^d)-d+2,\ldots,\Eseq(2^d)-1$ are impossible.
\end{theorem}

\begin{proof}
Induct on $d$.  Slice $S$ in a coordinate with at least two nonempty
fibers.  The defect-fiber inequality and the partition gap force the split to
be $(2^{d-1},2^{d-1})$ whenever the defect is above the stated gap.  The
induction hypothesis forces both projected slices to be $(d-1)$-boxes.  If
they coincide, their two lifts form a $d$-box.  If they differ, box rigidity
forces their intersection loss to be at least $2^{d-2}$, contradicting the
assumed defect.  The case $d=1$ is immediate.
\end{proof}

\begin{conjecture}[Defect continuity]\label{thm:continuity}
For a separately specified shape class, every integer defect in the proposed
interval from the spanning-tree value $R-1$ through $\Eseq(R-1)+1$ may be
achievable. The previous pendant-leaf induction is not a proof: a suitable
leaf need not exist, and the arithmetic overlap argument has counterexamples.
\end{conjecture}

The finite computation for $R=8$ found the defect values
$\{7,8,9,10,12\}$ in the tested connected-subset search. Theorem
\ref{thm:fracture} now explains the missing value $11$ for every subset; the
enumeration remains only a cross-check of the theorem and of the equality
orbit.

\begin{theorem}[Secondary fracture gap]\label{thm:secondary-gap}
Let $d\ge3$ and $|S|=2^d-1$.  If

\[
D(S)\ge \Eseq(2^d-2)+2,
\]

then $S$ is a $d$-box with one vertex removed.  Consequently, the values
$\Eseq(2^d-2)+2,\ldots,\Eseq(2^d-1)-1$ are impossible.
\end{theorem}

\begin{proof}
Slice in a coordinate and apply the multiway partition bound.  The only
partition that can reach the displayed range is
$(2^{d-1},2^{d-1}-1)$.  The larger projected slice is forced to be a
$(d-1)$-box by Theorem~\ref{thm:fracture}.  The complement identity for
$\Eseq$ gives a strict deficit for every nontrivial deviation of the smaller
slice from that box.  Hence it is the larger box with one vertex removed,
and lifting the two slices gives the result.
\end{proof}

\begin{conjecture}[Partial defect classification]\label{thm:classification}
Beyond the two gap theorems, a classification of achievable defect values may
consist of a low-defect interval together with isolated dense values. Any
general version must account for clique lines, embedding parameters, and the
distinction between defect and internal edge count.
\end{conjecture}

The topological funnel discussion in Section~\ref{sec:funnel} should therefore
be read as a program of conjectures and finite diagnostics except for the two
gap theorems above. The proven statements used by the paper include the defect
upper bound, the exact boundary identities, the power-of-two and near-full
fracture theorems, and the conditional restricted-boundary composition.
Continuity and the remaining equality classification stay open.
